\documentclass[11pt]{article}
\usepackage{amsmath,amssymb,amsthm,mathtools,booktabs,enumitem}
\usepackage[margin=1in]{geometry}
\title{Step 136: Response-Metric Normalization for the Incidence Operator}
\author{Six Birds / RH Membrane Working Notes}
\date{}

\newtheorem{theorem}{Theorem}
\newtheorem{lemma}{Lemma}
\newtheorem{definition}{Definition}
\newtheorem{proposition}{Proposition}
\newtheorem{remark}{Remark}

\begin{document}
\maketitle

\section{Purpose}
Step 135 showed that the correct high-divisibility tail is not the raw seed tail, but the response-space tail
\[
\tau_{\Omega,N}^{Z}
=
\left\|Z_y(I-P_{\Omega\le K})B_NG_{B,N}^{-1/2}\right\| .
\]
The immediate question is whether the Burnol--Müntz/semilocal response metric reduces the raw Boolean amplification
\[
\|Z_y\|=\varphi^k,
\qquad \varphi=\frac{1+\sqrt5}{2},
\]
or whether the actual seed tail must beat that full amplification.  Step 136 answers: the natural central-line Dirichlet response metric does reduce the amplification drastically, but does not eliminate it.

\section{Finite squarefree incidence model}
Let $P_y$ be a finite set of primes.  Identify squarefree divisors of $P_y$ with subsets $S\subseteq P_y$.  Define the incidence, or finite zeta-convolution, operator
\[
(Z_yb)(T)=\sum_{S\subseteq T}b(S).
\]
Locally, at one prime, the incidence matrix in the ordered basis $(0,1)$ is
\[
Z_p=\begin{pmatrix}1&0\\1&1\end{pmatrix}.
\]

\begin{definition}[Weighted coefficient metric]
For $\alpha\ge0$, define the squarefree coefficient metric
\[
\|b\|_{\alpha}^{2}
=
\sum_{S\subseteq P_y}\left(\prod_{p\in S}p^{-\alpha}\right)|b(S)|^2.
\]
The case $\alpha=0$ is raw Boolean counting.  The case $\alpha=1$ is the central-line Dirichlet coefficient metric corresponding to polynomials written as $\sum a_n n^{-1/2}\chi(n)$.
\end{definition}

\section{Local singular value}
In the normalized $\alpha$-metric, the local matrix is
\[
A_{p,\alpha}=W_p^{1/2}Z_pW_p^{-1/2}
=
\begin{pmatrix}1&0\\p^{-\alpha/2}&1\end{pmatrix}.
\]
Thus its largest singular value is
\[
\sigma_+(p,\alpha)
=
\frac{\sqrt{4+p^{-\alpha}}+p^{-\alpha/2}}{2}.
\]

\begin{theorem}[Metric-normalized incidence amplification]
For the finite incidence operator $Z_y$ on the squarefree cube, acting in the $\alpha$-weighted coefficient metric,
\[
\boxed{
\|Z_y\|_{\alpha\to\alpha}
=
\prod_{p\le y}\sigma_+(p,\alpha)
}
\]
where
\[
\sigma_+(p,\alpha)=\frac{\sqrt{4+p^{-\alpha}}+p^{-\alpha/2}}{2}.
\]
Consequently:
\[
\alpha=0:
\quad
\|Z_y\|=\varphi^{\pi(y)},
\]
\[
\alpha=1:
\quad
\log\|Z_y\|
=\frac12\sum_{p\le y}p^{-1/2}+O\!\left(\sum_{p\le y}p^{-1}\right)
=(1+o(1))\frac{\sqrt y}{\log y},
\]
\[
\alpha=2:
\quad
\log\|Z_y\|
=\frac12\sum_{p\le y}p^{-1}+O(1)
=\frac12\log\log y+O(1),
\]
while for $\alpha>2$ the product remains bounded.
\end{theorem}

\begin{proof}
The full squarefree incidence operator is the tensor product of the one-prime incidence matrices.  Under the weighted tensor-product metric, each one-prime normalized matrix is $A_{p,\alpha}$.  Tensor-product singular values multiply, giving the exact product formula.  The asymptotics follow from
\[
\log\sigma_+(p,\alpha)
=\frac12p^{-\alpha/2}+O(p^{-\alpha})
\]
and the prime number theorem estimates for sums over primes.
\end{proof}

\section{Interpretation for the membrane route}
The raw Boolean bound $\varphi^k$ is too pessimistic for the Dirichlet-source problem.  The natural $L$-function coefficient metric is $\alpha=1$, and in that metric the incidence amplification is only
\[
\exp\left((1+o(1))\frac{\sqrt y}{\log y}\right).
\]
This is a major reduction.  It is not a proof of tail smallness.  It says that a crude tail estimate of the form
\[
\|(I-P_{\Omega\le K})B_NG_B^{-1/2}\|\le e^{-cK}
\]
would imply
\[
\tau_{\\Omega,N}^{Z}
\lesssim
\exp\left((1+o(1))\frac{\sqrt y}{\log y}\right)e^{-cK}
\]
in the central-line Dirichlet metric.  Thus the seed tail must beat the weighted amplification, unless one uses the sharper blind-projected identity
\[
\Pi_{\mathcal B}Z_yB_N
=
\Pi_{\mathcal B}Z_y(I-P_{\Omega\le K})B_N,
\]
which may be substantially smaller than the full incidence norm.

\section{Response-pulled metric}
There is a formally perfect metric:
\[
\|b\|_{Z,H}^{2}=\|Z_yb\|_{H}^{2}.
\]
In that metric, $Z_y$ is an isometry by definition.  This is lawful only if the pulled metric is carried explicitly in the source and tail records.  It does not make the tail vanish; it moves the burden to the projection geometry of $P_{\Omega\le K}$ in the pulled response metric.

\section{Step 136 verdict}
\[
\boxed{
\text{The semilocal/Dirichlet response metric softens incidence amplification but does not remove it.}
}
\]
The next exact target is therefore not the full incidence norm, but the blind-projected norm
\[
\boxed{
\left\|\Pi_{\mathcal B}Z_y(I-P_{\Omega\le K})B_NG_{B,N}^{-1/2}\right\|.
}
\]
This is the response-metric-correct version of the GCD-log residual gate.

\section{Nonclaim boundary}
Step 136 does not prove RH, does not prove the actual Burnol/Müntz seed has small high-$\Omega$ tail, and does not prove that BPRZ supplies a full lower frame.  It only replaces the raw Boolean amplification estimate by the metric-correct amplification estimate, and isolates the next blind-projected tail gate.

\end{document}
